\scanpage{322}
\mspage{238}
% archive sheet number: 321

f et s. qu'elle \inserted{soit $0$-connexe, et} satisfasse $\deleted{\unsure{P}}\mathrm{PF}_{\unsure{\pi|\kappa}}2$ (cf page 81) ainsi que la variante suivante de $\mathrm{PF}_\pi1'$

$\mathrm{PF}_\pi1)$ Pour \deleted{\unsure{de|un}}tout objet $\deleted{\unsure{y}}x_0$ de $J$, $ {}_{x_0}\backslash J$ satisfait $\mathrm{PF}_\pi1'$, i.e. pour \unsure{tte|toute} famille de flèches $x_0\xrightarrow{u_\alpha}\deleted{\unsure{x}}y_\alpha$ $(\alpha\in A,\operatorname{card} A\leq\pi)$, existe une \deleted{\unsure{des}} flèche $x_0\xrightarrow{v}\deleted{\unsure{y}}z$ et des flèches $v_\alpha:y_\alpha\unsure{\mathbin{-}|\longrightarrow}z$, avec $v_\alpha u_\alpha=v$ $\forall\,\alpha\unsure{\in|+} A$.

Démonstration. Si $J$ est \unsure{grde|gde} devant $\pi$, alors \deleted{\unsure{alors}} $\mathrm{PF}_\pi1)$ résulte aussitôt de la conjonction de $\mathrm{PF}_\pi1'$ et $\mathrm{PF}_\pi2$, d'où le ``il faut''. Inversement, supposons que $J$ satisfasse $\mathrm{PF}_\pi1,\mathrm{PF}_\pi2$, et soit $0$-connexe, il faut prouver qu'elle satisfait $\mathrm{PF}_\pi1'$. Mais on sait déjà que les conditions en question, \inserted{qui} impliquent \deleted{que} \inserted{$\mathrm{PF}1$, $\mathrm{PF}2$, $0$-conn.,} \deleted{que} \inserted{impliquent que} $J$ est filtrante. Donc si on a une famille $(x_\alpha)_{\alpha\in A}$ \deleted{non vide}, $\operatorname{card} A\leq\pi$, choisissant \inserted{un} $x_0\in\mathrm{Ob}\,J$, on peut pour \unsure{tt|tout} $\alpha$ trouver $y_\alpha$ qui majore $x_0$ et $x_\alpha$, d'où des flèches

\diagram{\begin{tikzcd}[ampersand replacement=\&,row sep=tiny]
x_\alpha\arrow[dr,dotted]\&\\
\&y_\alpha,\\
x_0\arrow[ur]\&
\end{tikzcd}}

on peut \deleted{\unsure{leur}} appliquer $\mathrm{PF}_\pi1$, d'où des flèches $y_\alpha\longrightarrow z$, et en composant des $x_\alpha\mathbin{-}z$, qed.

\deleted{Corollaire} Prouvons maintenant le lemme. Il \deleted{\unsure{en}} suffit \inserted{(par le cor. précédent)} de prouver que $J$ satisfait à $\mathrm{PF}_\pi1$ et $\mathrm{PF}_\pi2$. Or $\mathrm{PF}_\pi1$ est immédiat, puisque $\forall\,x_0$ de $J$, $ {}_{x_0}\backslash J\xrightarrow{\sim}{}_{\unsure{y_0|\underline{x}_0}}\backslash I$, où $y_0$ est

